% The causal response theta_h = rho^h for two persistence values. Both start at
% the same impact; each period multiplies the response by rho, so by h = 7 one
% has all but vanished and the other is still near half its impact.
\documentclass[border=3pt]{standalone}
% ————— Shared preamble for the course figures —————
% Each figure is a standalone TikZ document. build.sh compiles it to a PDF (for
% the LaTeX build) and converts that to an SVG (for the web build).
%
% Nothing is drawn in pure black: the chrome sits at a mid grey that stays
% legible on white paper and on the site's dark theme alike. Data colours are
% the course palette, the same blue and orange used in the practicum.

\usepackage{amsmath}
\usepackage{tikz}
\usepackage{pgfplots}
\pgfplotsset{compat=1.18}
\usepgfplotslibrary{fillbetween,groupplots}
\usetikzlibrary{arrows.meta,positioning,calc,backgrounds}

\definecolor{nwBlue}{HTML}{2A78D6}
\definecolor{nwOrange}{HTML}{D9601F}
\definecolor{nwMut}{HTML}{898781}
\definecolor{nwInk}{HTML}{6B6965}

\pgfplotsset{
  nwaxis/.style={
    axis lines=left,
    axis line style={nwMut, line width=0.5pt},
    tick align=outside,
    tick style={nwMut, line width=0.5pt},
    label style={text=nwInk, font=\small},
    tick label style={text=nwMut, font=\footnotesize},
    legend cell align=left,
    legend style={draw=none, fill=none, font=\footnotesize, text=nwInk,
                  row sep=1pt, inner sep=1pt},
    every axis plot post/.append style={line join=round},
    clip mode=individual,
  },
}

\directlua{dofile("fig-l01-persistence-data.lua")}
\begin{document}
\begin{tikzpicture}
\begin{axis}[
  nwaxis,
  width=12cm, height=6.4cm,
  xmin=-0.4, xmax=12.6, ymin=0, ymax=1.1,
  xtick={0,1,...,12},
  ytick={0,0.25,0.5,0.75,1},
  xlabel={horizon $h$ (periods after the intervention)},
  ylabel={response $\theta_h$},
]

% guide at h = 7
\addplot[nwMut, densely dotted, line width=0.6pt]
  coordinates {(7,0) (7,1.02)};

\addplot[nwOrange, line width=1.2pt, mark=*, mark size=1.7pt,
         mark options={fill=nwOrange, draw=nwOrange}] coordinates {\lpPsLo};
\addplot[nwBlue, line width=1.3pt, mark=*, mark size=1.7pt,
         mark options={fill=nwBlue, draw=nwBlue}] coordinates {\lpPsHi};

% direct labels
\node[text=nwBlue, font=\footnotesize, anchor=south west]
  at (axis cs:3.2,0.8) {$\rho=0.9$};
\node[text=nwOrange, font=\footnotesize, anchor=south west]
  at (axis cs:2.15,0.3) {$\rho=0.5$};

% same impact
\node[text=nwInk, font=\footnotesize, anchor=west]
  at (axis cs:0.25,1.035) {same impact, $\theta_0=1$};

% the h = 7 comparison
\node[text=nwBlue, font=\footnotesize, anchor=west]
  at (axis cs:7.2,\lpPsHiSeven+0.06) {$0.9^7=\lpPsHiSevenLab$};
\node[text=nwOrange, font=\footnotesize, anchor=south west]
  at (axis cs:7.2,\lpPsLoSeven+0.03) {$0.5^7=\lpPsLoSevenLab$};

\end{axis}
\end{tikzpicture}
\end{document}
